% Pursuing Stacks (cotes 134-2 to 134-8): archive cote 134-7, site page 1 (PDF page 2), typescript page 443.
% typescript. Draft compared: ps6.tex:1-58.
% AI-transcribed by Claude Opus 5.5 (claude-opus-5-5, adaptive thinking, effort high), Message Batches API, batch ps3 (msgbatch_01XdFmuVVBAcgc7EFxdtCqsj), 2026-11-12; not yet checked by eye.
\scanpage{134-7/1}
\margin{\uline{VI Schematiza-tion}}

\mspage{443}

24.8.

110.\quad I pondered some more about homotopy types over a grou\unsure{nd r}ing $k$,
% NOTE: a rust stain runs down the right side of the page and hides letters at many line ends
just enough to become familiar again with the idea\hand{,}and more or less con-%
vinced that there should exist such a thing, which \xout{amounts} shoul\unsure{d} amount\hand{,}
kind of\hand{,} to putting a ``continuous'' structure \over{\unsure{s}}{(}namely the very rich struc-%
ture of a s\unsure{\hand{c}}heme) upon something usually visualized as something ``discrete''
- namely a homotopy type. The basic analogy here is free $\underline{Z}$-modules $M$
of finite type - a typical case of a ``\uline{discrete}'' structure. It gives
rise though to a \uline{vector bundle} $W(M)$ over the absolute base $S_o=\mathrm{Spec}(\underline{Z})$,
whose $\underline{Z}$-module of sections is $M$ , and the functor
\[ M \mapsto W(M) \]
from free $\underline{Z}$-modules of finite type to vector bundles over $S_o$ is fully
faithful. When $M$ is an arbitrary $\underline{Z}$-module i.e. an abelian group, $W(M)$
still makes a sense, namely as a functor
\[ k \mapsto M\otimes_{\underline{Z}}k \]
on the category of all commutative $\underline{Z}$-algebras (i.e. just commutative
rings); it is no longer representable by a scheme over $S_o$ (except \deleted{when}
precisely \inserted{when} $M$ is free of finite type), but it is very close \xout{\illegible{1}}\inserted[typed]{still}, intui-%
tively and technically too, to a usual vector bundle (the ``vect\unsure{or}'' struc-%
ture coming from the $k$-module structure on $W(M)(k) = M\otimes_{\underline{Z}}k$ ). Again,
the functor $M\mapsto W(M)$ from Ab to the category of ``generalized \unsure{ve}ctor
bundles'' over $S_o$ is fully faithful. Working with semisimplic\unsure{ia}l
$\underline{Z}$-modules (say) rather than \unsure{\over{j}{j}}ust $\underline{Z}$-modules, and more specifical\unsure{ly}
with those corresponding to $K(\hand{\pi},n)$ types, and using Postnikov ``\unsure{dé}vissa-%
ge'' of a general homotopy type, one may hope to ``represent'' th\unsure{is} type,
in a more or less canonical way in terms of the successive semisi\unsure{m}plicial
P\hand{\u{o}}stnikov fibrations, by a semi-simplicial object in the topo\unsure{s} (\unsure{s}ay) of
all functors from $\underline{Z}$-algebras to sets which are ``sheaves'' for a s\unsure{u}itable
site structure on the dual category (namely the category of affi\unsure{ne}
sch\unsure{em|mm}es over $\underline{Z}$ ) - the so-called ``flat'' topology seems \xout{to} OK. (NB To
eliminate logical difficulties, we may have to restrict somewhat the
rings $k$ used as arguments, for instance take them to be of fini\unsure{te} type
over $\underline{Z}$ - never mind such technicalities !). This approach \xout{is} ma\unsure{y}
possibly work, when restricting to \xout{simply} 1-connected homotopy \unsure{t}ypes,
or at any rate to the case when the fundamental groups of the con\unsure{n}ected
\margin{*} components are abelian. If we wish a ``schemati\unsure{z}ation'' of arbitrary
homotopy types, we may think of going about it by keeping the f\unsure{u}ndamen-%
tal group or groupoid ``discrete'', and ``schematize'' the P\hand{\u{o}}stni\unsure{ko}v

\margin{1}
